\documentclass[12pt,a4paper]{article}

\usepackage[utf8]{inputenc}
\usepackage[brazilian]{babel}
\usepackage{amsmath,amssymb}
\usepackage{graphicx}
\usepackage{float}
\usepackage{geometry}
\usepackage{booktabs}
\usepackage{hyperref}
\usepackage{listings}
\usepackage{xcolor}

\geometry{margin=2.5cm}
\hypersetup{colorlinks=true, linkcolor=black, urlcolor=blue}
\graphicspath{{resultados/figuras/}}

\lstset{
    basicstyle=\ttfamily\small,
    breaklines=true,
    frame=single,
    numbers=left,
    numberstyle=\tiny,
    keywordstyle=\color{blue!70!black},
    commentstyle=\color{green!40!black}
}

\title{Terceira Avaliação de Métodos Numéricos}
\author{Gustavo Bittar Gonçalves}
\date{24 de junho de 2026}

\begin{document}
\maketitle

\section{Condução térmica bidimensional transiente}

O problema físico corresponde à condução de calor em uma placa bidimensional com propriedades constantes e sem geração interna. O modelo diferencial é
\begin{equation}
    \rho c_p \frac{\partial T}{\partial t}
    =
    k\left(
        \frac{\partial^2 T}{\partial x^2}
        +
        \frac{\partial^2 T}{\partial y^2}
    \right),
    \qquad
    0<x<L_x,\quad 0<y<L_y,\quad t>0.
\end{equation}
Equivalentemente,
\begin{equation}
    \frac{\partial T}{\partial t}
    =
    \alpha\left(
        \frac{\partial^2 T}{\partial x^2}
        +
        \frac{\partial^2 T}{\partial y^2}
    \right),
    \qquad
    \alpha=\frac{k}{\rho c_p}.
\end{equation}

As condições de contorno gerais são de Dirichlet:
\begin{align}
    T(0,y,t)&=T_E, & T(L_x,y,t)&=T_D,\\
    T(x,0,t)&=T_S, & T(x,L_y,t)&=T_N,
\end{align}
e a condição inicial usada no domínio interno é
\begin{equation}
    T(x,y,0)=300\ \mathrm{K}.
\end{equation}

Para o caso numérico da avaliação:
\begin{align}
    T(0,y,t)&=400\ \mathrm{K},&
    T(L_x,y,t)&=500\ \mathrm{K},\\
    T(x,0,t)&=600\ \mathrm{K},&
    T(x,L_y,t)&=200\ \mathrm{K}.
\end{align}

Foram usadas as propriedades do alumínio puro
\[
    k=237\ \mathrm{W/(m\,K)},\qquad
    \rho=2700\ \mathrm{kg/m^3},\qquad
    c_p=900\ \mathrm{J/(kg\,K)},
\]
logo
\[
    \alpha=9{,}7531\times 10^{-5}\ \mathrm{m^2/s}.
\]
A placa possui $L_x=L_y=1{,}0\ \mathrm{m}$ e espessura de $5{,}0\ \mathrm{mm}$. Como não há termo fonte nem perda superficial, a espessura não aparece explicitamente na equação bidimensional resolvida.

\subsection{Discretização implícita}

Aplicando diferenças finitas centrais de segunda ordem no espaço e avanço temporal implícito, tem-se
\begin{equation}
    \frac{T_{i,j}^{n+1}-T_{i,j}^{n}}{\Delta t}
    =
    \alpha
    \left[
        \frac{T_{i+1,j}^{n+1}-2T_{i,j}^{n+1}+T_{i-1,j}^{n+1}}{\Delta x^2}
        +
        \frac{T_{i,j+1}^{n+1}-2T_{i,j}^{n+1}+T_{i,j-1}^{n+1}}{\Delta y^2}
    \right].
\end{equation}
Definindo
\[
    r_x=\frac{\alpha\Delta t}{\Delta x^2},
    \qquad
    r_y=\frac{\alpha\Delta t}{\Delta y^2},
\]
obtém-se, para cada ponto interno,
\begin{equation}
    (1+2r_x+2r_y)T_{i,j}^{n+1}
    -r_x(T_{i+1,j}^{n+1}+T_{i-1,j}^{n+1})
    -r_y(T_{i,j+1}^{n+1}+T_{i,j-1}^{n+1})
    =
    T_{i,j}^{n}.
\end{equation}
O passo de tempo foi calculado por
\begin{equation}
    \Delta t=C\frac{\min(\Delta x,\Delta y)^2}{\alpha}.
\end{equation}

\subsection{Algoritmo de Gauss-Seidel}

Isolando $T_{i,j}^{n+1}$, a iteração de Gauss-Seidel fica
\begin{equation}
    T_{i,j}^{(m+1)}
    =
    \frac{
        T_{i,j}^{n}
        +r_x\left(T_{i-1,j}^{(m+1)}+T_{i+1,j}^{(m)}\right)
        +r_y\left(T_{i,j-1}^{(m+1)}+T_{i,j+1}^{(m)}\right)
    }{1+2r_x+2r_y}.
\end{equation}
No programa foi usada a forma Gauss-Seidel vermelho-preto vetorizada, que preserva a dependência atualizada entre vizinhos e reduz o custo computacional em Python.

\begin{enumerate}
    \item Inicializar $T^0$ com 300 K no interior e impor as temperaturas prescritas no contorno.
    \item Para cada passo de tempo, calcular $r_x$ e $r_y$.
    \item Repetir Gauss-Seidel até que $\max|T^{(m+1)}-T^{(m)}|<10^{-8}$.
    \item Atualizar o campo temporal e salvar os instantes pedidos.
\end{enumerate}

\begin{table}[H]
\centering
\caption{Parâmetros efetivos da simulação com $N_x=N_y=10$.}
\begin{tabular}{ccccc}
\toprule
$C$ & $\Delta t$ (s) & $r_x$ & $r_y$ & iterações médias GS \\
\midrule
0,25 & 31,578947 & 0,249474 & 0,249474 & 7,03 \\
0,50 & 63,157895 & 0,498947 & 0,498947 & 10,92 \\
1,00 & 126,315790 & 0,997895 & 0,997895 & 18,08 \\
2,00 & 252,631580 & 1,995790 & 1,995790 & 30,82 \\
\bottomrule
\end{tabular}
\end{table}

\subsection{Campos de temperatura}

\begin{figure}[H]
\centering
\includegraphics[width=\textwidth]{calor_campos_C0.25.png}
\caption{Campos de temperatura para $C=0{,}25$.}
\end{figure}

\begin{figure}[H]
\centering
\includegraphics[width=\textwidth]{calor_campos_C0.5.png}
\caption{Campos de temperatura para $C=0{,}50$.}
\end{figure}

\begin{figure}[H]
\centering
\includegraphics[width=\textwidth]{calor_campos_C1.png}
\caption{Campos de temperatura para $C=1{,}0$.}
\end{figure}

\begin{figure}[H]
\centering
\includegraphics[width=\textwidth]{calor_campos_C2.png}
\caption{Campos de temperatura para $C=2{,}0$.}
\end{figure}

\subsection{Perfis em $t=4$ h}

Para os perfis de malha foi usado $C=1{,}0$, mantendo o método implícito e comparando $N_x=N_y=10$, $20$ e $40$.

\begin{figure}[H]
\centering
\includegraphics[width=0.86\textwidth]{calor_perfil_x_meio.png}
\caption{Perfil em $x=L_x/2$.}
\end{figure}

\begin{figure}[H]
\centering
\includegraphics[width=0.86\textwidth]{calor_perfil_x_quarto.png}
\caption{Perfil em $x=L_x/4$.}
\end{figure}

\begin{figure}[H]
\centering
\includegraphics[width=0.86\textwidth]{calor_perfil_y_meio.png}
\caption{Perfil em $y=L_y/2$.}
\end{figure}

\begin{figure}[H]
\centering
\includegraphics[width=0.86\textwidth]{calor_perfil_y_3quartos.png}
\caption{Perfil em $y=3L_y/4$.}
\end{figure}

\subsection{Evolução temporal}

\begin{figure}[H]
\centering
\includegraphics[width=0.86\textwidth]{calor_evolucao_centro.png}
\caption{Evolução temporal da temperatura em $x=L_x/2$ e $y=L_y/2$.}
\end{figure}

\begin{figure}[H]
\centering
\includegraphics[width=0.86\textwidth]{calor_evolucao_quarto.png}
\caption{Evolução temporal da temperatura em $x=L_x/4$ e $y=L_y/4$.}
\end{figure}

\section{Equação da onda unidimensional}

A equação da onda é
\begin{equation}
    \frac{\partial^2 y}{\partial t^2}
    =
    c^2\frac{\partial^2 y}{\partial x^2},
    \qquad 0\le x\le L,\quad 0\le t\le 2\ \mathrm{s}.
\end{equation}
Foram usados $L=1\ \mathrm{m}$, $c=1\ \mathrm{m/s}$ e $A=0{,}01\ \mathrm{m}$, com
\begin{equation}
    y(0,t)=y(L,t)=0,
    \qquad
    y(x,0)=A\sin\left(\frac{\pi x}{L}\right),
    \qquad
    v(x,0)=\frac{\partial y}{\partial t}(x,0)=0.
\end{equation}

Definindo $v=\partial y/\partial t$, o sistema de primeira ordem é
\begin{equation}
    \frac{\partial y}{\partial t}=v,
    \qquad
    \frac{\partial v}{\partial t}=c^2\frac{\partial^2 y}{\partial x^2}.
\end{equation}
No espaço,
\begin{equation}
    \frac{d y_i}{dt}=v_i,
    \qquad
    \frac{d v_i}{dt}=c^2\frac{y_{i+1}-2y_i+y_{i-1}}{\Delta x^2}.
\end{equation}

\subsection{Euler explícito}

\begin{align}
    y_i^{n+1} &= y_i^n+\Delta t\,v_i^n,\\
    v_i^{n+1} &= v_i^n+\Delta t\,c^2
    \frac{y_{i+1}^n-2y_i^n+y_{i-1}^n}{\Delta x^2}.
\end{align}
Esse método, aplicado diretamente ao sistema de primeira ordem da onda, é instável para o oscilador sem amortecimento. Por isso os resultados de erro crescem rapidamente.

\subsection{Runge-Kutta explícito de quarta ordem}

Se $\mathbf{u}=[\mathbf{y},\mathbf{v}]^T$ e $\mathbf{u}'=F(\mathbf{u})$, então
\begin{align}
    \mathbf{k}_1 &= F(\mathbf{u}^n),\\
    \mathbf{k}_2 &= F(\mathbf{u}^n+\tfrac{\Delta t}{2}\mathbf{k}_1),\\
    \mathbf{k}_3 &= F(\mathbf{u}^n+\tfrac{\Delta t}{2}\mathbf{k}_2),\\
    \mathbf{k}_4 &= F(\mathbf{u}^n+\Delta t\mathbf{k}_3),\\
    \mathbf{u}^{n+1} &= \mathbf{u}^n+\frac{\Delta t}{6}
    (\mathbf{k}_1+2\mathbf{k}_2+2\mathbf{k}_3+\mathbf{k}_4).
\end{align}

\subsection{Euler implícito}

Com $\mathbf{D}_{xx}$ representando a matriz de segunda derivada nos pontos internos,
\begin{align}
    \mathbf{y}^{n+1}-\Delta t\,\mathbf{v}^{n+1}&=\mathbf{y}^n,\\
    -\Delta t\,c^2\mathbf{D}_{xx}\mathbf{y}^{n+1}+\mathbf{v}^{n+1}&=\mathbf{v}^n.
\end{align}
Logo, o sistema linear resolvido a cada passo é
\begin{equation}
    \begin{bmatrix}
        \mathbf{I} & -\Delta t\,\mathbf{I}\\
        -\Delta t\,c^2\mathbf{D}_{xx} & \mathbf{I}
    \end{bmatrix}
    \begin{bmatrix}
        \mathbf{y}^{n+1}\\
        \mathbf{v}^{n+1}
    \end{bmatrix}
    =
    \begin{bmatrix}
        \mathbf{y}^{n}\\
        \mathbf{v}^{n}
    \end{bmatrix}.
\end{equation}
A matriz $\mathbf{D}_{xx}$ é tridiagonal:
\begin{equation}
    \mathbf{D}_{xx}
    =
    \frac{1}{\Delta x^2}
    \begin{bmatrix}
        -2 & 1 & 0 & \cdots & 0\\
        1 & -2 & 1 & \cdots & 0\\
        0 & 1 & -2 & \cdots & 0\\
        \vdots & \vdots & \vdots & \ddots & 1\\
        0 & 0 & 0 & 1 & -2
    \end{bmatrix}.
\end{equation}

\subsection{Solução analítica e erro}

A solução analítica usada para a comparação é
\begin{equation}
    y(x,t)=A\sin\left(\frac{\pi x}{L}\right)\cos\left(\frac{\pi c t}{L}\right).
\end{equation}
O erro discreto foi calculado por
\begin{equation}
    E(t)=\sqrt{\frac{1}{N}\sum_{i=1}^{N}\left(y_i^{\mathrm{num}}-y_i^{\mathrm{exata}}\right)^2}.
\end{equation}

\subsection{Resultados da onda}

As simulações usaram $N_x=101$ e $\Delta t=0{,}004$ s.

\begin{figure}[H]
\centering
\includegraphics[width=0.86\textwidth]{onda_deslocamento_euler_explicito.png}
\caption{Deslocamento $y(x,t)$ por Euler explícito.}
\end{figure}

\begin{figure}[H]
\centering
\includegraphics[width=0.86\textwidth]{onda_deslocamento_rk4.png}
\caption{Deslocamento $y(x,t)$ por Runge-Kutta de quarta ordem.}
\end{figure}

\begin{figure}[H]
\centering
\includegraphics[width=0.86\textwidth]{onda_deslocamento_euler_implicito.png}
\caption{Deslocamento $y(x,t)$ por Euler implícito.}
\end{figure}

\begin{figure}[H]
\centering
\includegraphics[width=0.86\textwidth]{onda_comparacao_t2.png}
\caption{Comparação dos métodos em $t=2$ s.}
\end{figure}

\begin{figure}[H]
\centering
\includegraphics[width=0.86\textwidth]{onda_evolucao_y_centro.png}
\caption{Evolução temporal de $y(L/2,t)$.}
\end{figure}

\begin{figure}[H]
\centering
\includegraphics[width=0.86\textwidth]{onda_evolucao_v_centro.png}
\caption{Evolução temporal de $v(L/2,t)$.}
\end{figure}

\begin{figure}[H]
\centering
\includegraphics[width=0.86\textwidth]{onda_erro_l2.png}
\caption{Evolução temporal do erro $L_2$.}
\end{figure}

\section{Programas computacionais}

O código completo está no arquivo \texttt{main.py}. Ele executa todas as simulações e grava automaticamente as figuras usadas neste relatório. Para atender ao pedido de três programas computacionais da questão 9, também foram criados arquivos de execução separados para cada método da onda:
\[
\texttt{onda\_euler\_explicito.py},\qquad
\texttt{onda\_rk4.py},\qquad
\texttt{onda\_euler\_implicito.py}.
\]
Cada um desses arquivos chama a mesma implementação numérica validada em \texttt{main.py}, evitando divergência entre o relatório e os programas.

\begin{lstlisting}[language=Python]
def run_all() -> None:
    ensure_dirs()
    heat_cfg = HeatConfig()
    heat_times = [0.0, 1800.0, 3600.0, 7200.0, 10800.0, 14400.0]
    heat_c_results = {}
    for c_value in [0.25, 0.50, 1.0, 2.0]:
        result = simulate_heat(10, 10, c_value, heat_times, heat_cfg)
        heat_c_results[c_value] = result
        plot_heat_snapshots(result, c_value)

    heat_mesh_results = {}
    for n in [10, 20, 40]:
        heat_mesh_results[n] = simulate_heat(n, n, 1.0, [0.0, 14400.0], heat_cfg)
    plot_heat_mesh_profiles(heat_mesh_results, heat_cfg)
    plot_heat_temporal(heat_mesh_results)

    wave_cfg = WaveConfig()
    wave_results = {
        "euler_explicito": simulate_wave("euler_explicito", wave_cfg),
        "rk4": simulate_wave("rk4", wave_cfg),
        "euler_implicito": simulate_wave("euler_implicito", wave_cfg),
    }
    plot_wave_results(wave_results, wave_cfg)
    save_summary(heat_c_results, wave_results)
\end{lstlisting}

\section{Conclusões}

O método implícito para condução térmica produziu campos suaves e convergentes para todos os coeficientes $C$ testados. O aumento de $C$ elevou o número médio de iterações de Gauss-Seidel por passo, como esperado, pois o acoplamento espacial do sistema linear se torna mais forte.

Na equação da onda, o método RK4 acompanhou muito bem a solução analítica no intervalo simulado. O Euler implícito permaneceu estável, mas introduziu amortecimento numérico. O Euler explícito aplicado diretamente ao sistema de primeira ordem apresentou crescimento instável, refletido no erro $L_2$.

\section{Referência}

INCROPERA, F. P.; DEWITT, D. P.; BERGMAN, T. L.; LAVINE, A. S. \textit{Fundamentos de transferência de calor e de massa}. LTC.

\end{document}
